\documentclass{sbc2019}%

\usepackage{graphicx}
\usepackage[utf8]{inputenc}
%\usepackage[brazil]{babel}
\usepackage[misc,geometry]{ifsym}
\usepackage{fontspec}
\usepackage{fontawesome}
\usepackage{academicons}
\usepackage{color}
\usepackage{hyperref}
\usepackage{aas_macros}
\usepackage[bottom]{footmisc}
\usepackage{booktabs}
\usepackage{multirow}

\definecolor{orcidlogo}{rgb}{0.37,0.48,0.13}
\definecolor{unilogo}{rgb}{0.16, 0.26, 0.58}
\definecolor{maillogo}{rgb}{0.58, 0.16, 0.26}
\definecolor{darkblue}{rgb}{0.0,0.0,0.0}
\hypersetup{colorlinks,breaklinks,
            linkcolor=darkblue,urlcolor=darkblue,
            anchorcolor=darkblue,citecolor=darkblue}

%%%%%%% IMPORTANT: We disable hyperlinks by default with this line, to avoid the error "\pdfendlink ended up in different nesting level" while writing.
%\hypersetup{draft}

\jid{JSERD}
\jtitle{Journal of Software Engineering Research and Development, 2026, }
\doi{10.5753/jserd.2026.xxx}
\copyrightstatement{This work is licensed under a Creative Commons Attribution 4.0 International License.}
\jyear{2026}


\title[Architectural Simplification and Internal Software Quality]{Architectural Simplification and Internal Software Quality in a Unity-Based Serious Game: A Comparative Case Study}

%% NEEDS SOURCE: the definitive author list and the ORCID identifiers for this
%% paper have not been fixed yet. The two authors below are the ones documented
%% in the PIBIC 2025--2026 research plan and final report.
\author[Rezende and Rocha 2026]{
\affil{\textbf{Jo\~ao Guilherme Benjamin Alves de Rezende}~~[~\textbf{Universidade Federal do Agreste de Pernambuco~}|\href{mailto:benj7100@gmail.com}{~\textbf{\textit{benj7100@gmail.com}}} ]}

\affil{\textbf{Rodrigo Gusm\~ao de Carvalho Rocha}~~[~\textbf{Universidade Federal do Agreste de Pernambuco~}|\href{mailto:rodrigo.rocha@ufape.edu.br}{~\textbf{\textit{rodrigo.rocha@ufape.edu.br}}} ]}
}

\begin{document}

\begin{frontmatter}
\maketitle

\begin{abstract}
\textbf{Abstract}

%% ABSTRACT -- written last, see Section ordering in the writing protocol.
Architectural refactoring is commonly reported as an activity that adds structure to a system. Less is known about the opposite direction: what happens to internal software quality indicators when abstractions that do not carry a concrete responsibility are removed from a working system. This paper reports a comparative case study of responsibility-driven architectural simplification applied to Molz, a Unity-based serious game for organic chemistry. Three snapshots of the same system were measured with a single frozen instrument: the pre-intervention baseline C, an intermediate control state D0 obtained by removing dead code only, and the simplified state D. Separating the control from the intervention was necessary because the code removed in C$\rightarrow$D0 had a higher mean maintainability index (80.3 versus 72.5) and lower mean efferent coupling (4.50 versus 11.4) than the system as a whole, so its removal moved every normalized mean unfavourably. Between D0 and D the system lost five classes and gained two interfaces and six members, while the aggregate indicators barely moved: maintainability index 71.7 to 71.8, cyclomatic complexity 2.39 to 2.38, and mean efferent coupling 12.8 to 13.4. A matched-population analysis over the 442 members present in both states inverted the sign of the first two: maintainability fell by 0.06 and complexity rose by 0.01, showing that the aggregate movements were composition effects rather than changes in surviving code. The set of types referencing the engine was identical in both states (23 types), while its percentage rose from 35.4\% to 38.3\% through denominator change alone. Functional behaviour was preserved under 118 characterization and architecture-invariant tests, a clean engine import, and manual full-flow verification. Two a-priori simplification hypotheses were rejected by implementation evidence rather than confirmed. The study contributes a protocol for responsibility-driven simplification with recorded retention decisions, a control design that separates dead-code removal from architectural intervention, and evidence that structural simplification does not translate into uniform movement of aggregate quality indicators.
\end{abstract}

\begin{keywords}
Software architecture, Architectural refactoring, Internal software quality, Maintainability, Coupling, Game software engineering, Unity, Case study
\end{keywords}

\end{frontmatter}


%% =========================================================================
\section{Introduction}
%% =========================================================================

Refactoring is defined as a change to the internal structure of software that
makes it easier to understand and cheaper to modify without altering its
observable behaviour \cite{fowler2018,mens2004}. In empirical studies the
activity is most often observed in one direction: structure is added. Classes
are extracted, interfaces are introduced, responsibilities are distributed
across new components. The catalogues of code smells contain the opposite case
as well --- Speculative Generality, Lazy Class, Middle Man, Dead Code
\cite{fowler2018} --- but the evidence base for removing structure is thinner
than the evidence base for adding it, and thinner still at the level of
architecture rather than of individual classes.

This matters because the two directions do not behave symmetrically under
measurement. Aggregate indicators such as the maintainability index
\cite{coleman1994}, cyclomatic complexity \cite{mccabe1976} and efferent
coupling \cite{martin2003} are means over a population of members or types.
When structure is added, the population grows; when structure is removed, the
population shrinks. A mean can therefore move without any surviving component
having changed at all. Studies that report only aggregate deltas cannot
distinguish the two, and the literature already reports that refactoring
produces mixed rather than uniform metric outcomes
\cite{szoke2017,bavota2015,palomba2018}.

Game software adds a further dimension. Games are built on engines that impose
their own composition model, lifecycle and serialization mechanics
\cite{aleem2016,politowski2025}, and engine-specific structures are a
recognized source of maintainability problems
\cite{nardone2023,borrelli2020}. In an engine-based system, a component's
dependency on the engine is not a free design variable: it is partly fixed by
the way the engine instantiates and wires objects. Whether architectural
simplification can reduce that dependency, and what happens to internal quality
indicators when it is attempted, is an open empirical question.

This paper reports a comparative case study of \emph{responsibility-driven
architectural simplification} in Molz, a Unity-based serious game for the
solubility classification of organic compounds. The guiding principle of the
intervention was not to impose a target architectural style, but to reassess
each existing abstraction against the responsibility it actually performed in
the running system, and to retain only those carrying a concrete one.

The study is organized around three snapshots of the same system. State C is
the pre-intervention baseline. State D0 is an intermediate control obtained by
removing dead code and nothing else. State D is the result of eight
architectural stages applied to D0. The control checkpoint was introduced after
an audit found that a substantial part of the baseline was unreachable; without
it, the metric consequences of deleting unreachable code would have been
credited to the architectural intervention.

\medskip\noindent\textbf{Objective.} To investigate how responsibility-driven
architectural simplification of a functionally preserved Unity-based serious
game changed its structural organization and internal software quality
indicators, distinguishing the effect of dead-code removal from that of the
subsequent architectural intervention.

\medskip\noindent\textbf{RQ1.} How did the system's architectural organization
change when existing abstractions were reassessed according to their concrete
responsibilities and redundant architectural elements were simplified?

\medskip\noindent\textbf{RQ2.} How did internal software quality indicators ---
maintainability, cyclomatic complexity, structural coupling, engine dependency,
and structural size --- differ between the pre-intervention state D0 and the
simplified state D, and how do those differences relate to changes in component
population?

\medskip\noindent\textbf{RQ3.} How were changes in internal software quality
distributed across components removed as dead code, directly refactored,
indirectly affected, structurally preserved, and newly introduced by the
architectural intervention?

\medskip

The contributions of the paper are four. First, a protocol for
responsibility-driven architectural simplification in which retention decisions
are recorded with the same weight as removal decisions, and in which an a-priori
hypothesis can be rejected by implementation evidence. Second, a C/D0/D
comparative design that quantitatively separates dead-code removal from
architectural intervention under verified functional preservation and a frozen
measurement instrument. Third, evidence that structural simplification did not
produce uniform movement of aggregate indicators, together with an explicit
decomposition of within-component change against composition effect --- including
two indicators whose sign inverts once the population is controlled. Fourth, a
component-level analysis showing that the measured structural changes were
confined to the components classified as directly or indirectly affected, while
the source-preserved group remained metric-identical.

The remainder of the paper is organized as follows. Section~\ref{sec:related}
positions the study against work on internal quality measurement, architectural
refactoring and game software architecture. Section~\ref{sec:design} describes
the case study design. Section~\ref{sec:results} reports the results for the
three research questions. Section~\ref{sec:discussion} interprets them.
Section~\ref{sec:threats} discusses threats to validity, and
Section~\ref{sec:conclusion} concludes.


%% =========================================================================
\section{Background and Related Work}\label{sec:related}
%% =========================================================================

\subsection{Internal Software Quality and Architecture}

Maintainability is one of the product-quality characteristics of ISO/IEC 25010
\cite{iso25010}, but the standard does not prescribe how to estimate it from
source code. Practice has settled on a small set of structural proxies.

The \emph{maintainability index} (MI) combines Halstead volume, cyclomatic
complexity and lines of code into a single composite score
\cite{coleman1994}. Its convenience is also its principal weakness: the
composite was calibrated on a specific corpus, and the relationship between the
score and the maintenance effort actually experienced by developers is
indirect. Systematic reviews of maintainability prediction report wide
variation in which metrics are used and how they are validated
\cite{riaz2009,ardito2020}. An influential alternative line of work argues that
maintainability is better assessed by mapping ISO characteristics onto several
separately reported source-code properties --- volume, unit size, unit
complexity, duplication, unit testing --- rather than collapsing them into one
number \cite{heitlager2007}, and by benchmarking those properties against a
population of comparable systems \cite{baggen2012}.

Cyclomatic complexity \cite{mccabe1976} has attracted a parallel critique.
Shepperd \cite{shepperd1988} argued that the measure rests on weak theoretical
foundations and that, for a large class of software, it is largely a proxy for
size. Since MI incorporates both cyclomatic complexity and lines of code, the
two critiques compound.

Coupling has a firmer conceptual basis. Chidamber and Kemerer
\cite{chidamber1994} defined coupling between objects as the number of other
classes a class depends upon, and Martin \cite{martin2003} defined efferent
coupling at package level together with the dependency inversion principle,
under which modules should depend on abstractions rather than on concrete
implementations. Efferent coupling counts outgoing dependencies; it says how
many other types a type must know about, not whether knowing about them is
justified.

We therefore treat MI, cyclomatic complexity and efferent coupling as
\emph{internal software quality indicators} and not as measurements of
maintainability itself. This distinction is not rhetorical: as
Section~\ref{sec:results} shows, two of these indicators move in one direction
in aggregate and in the opposite direction once the population is held fixed.

\subsection{Architectural Refactoring and Simplification}

Surveys of refactoring \cite{mens2004} and the catalogue of smells and
corresponding transformations \cite{fowler2018} cover both the addition and the
removal of structure. Empirically, however, the reported outcomes of refactoring
on internal quality attributes are mixed rather than uniformly positive.
Sz\H{o}ke et al.\ \cite{szoke2017}, studying large-scale industrial systems,
found that a single refactoring had only a small and sometimes negative effect
on maintainability, while the accumulated process had a beneficial one. Bavota
et al.\ \cite{bavota2015} found that refactoring operations are not a direct
consequence of poor metric profiles or of the presence of smells. Palomba et
al.\ \cite{palomba2018} found that the impact of smells on maintainability
varies substantially by smell type and by system.

Removal of structure has been studied most systematically in the specific case
of dead code. Romano et al.\ \cite{romano2020} report a multi-study
investigation showing that dead code is widespread and that developers perceive
it as harmful to comprehensibility, and Eder et al.\ \cite{eder2012} measured
unused code in an industrial system and its relation to maintenance effort.
This literature establishes that removing unreachable code is worthwhile; it
does not report what removing it does to the aggregate metric profile of the
system that remains --- which is precisely the confound the control state D0 in
this study is designed to isolate.

At the architectural level, the relevant body of work is architectural
technical debt. Verdecchia et al.\ \cite{verdecchia2020,verdecchia2021} built
and evaluated a grounded theory of architectural technical debt, characterizing
it as design decisions that were suitable when taken but later hinder progress
--- a description that fits abstractions whose justification has since
disappeared. Kruchten et al.\ \cite{kruchten2012} discuss the metaphor and its
limits, and Li et al.\ \cite{li2022} map the literature on architecture erosion
and its symptoms. This work is largely diagnostic: it characterizes and detects
degradation. Studies that apply a simplification intervention to a working
system and then measure what happened, with a control for confounding
deletions, are considerably rarer.

\subsection{Game Software Architecture and Engine Dependency}

Game development has been characterized as a distinct software engineering
context, with its own process life cycle and its own recurring difficulties
\cite{aleem2016}. Engine architecture has recently received direct attention:
Politowski et al.\ \cite{politowski2025} recovered and compared the
architectures of ten open-source engines, identifying excessive coupling as a
recurrent issue and showing in a controlled experiment that explicit
architectural models let developers perform understanding and impact-analysis
tasks faster and more correctly. Ullmann et al.\ \cite{ullmann2023} visualized
subsystem coupling patterns across engines and found the low-level renderer,
platform-independence layer and resource manager to be frequently coupled to
the engine core.

Closer to the present case, Nardone et al.\ \cite{nardone2023} derived a
catalogue of 28 video-game bad smells from developer discussions, and Borrelli
et al.\ \cite{borrelli2020} implemented static detection of seven Unity-specific
smells covering performance, maintainability and incorrect behaviour. Both
establish that engine-specific structures are a legitimate object of internal
quality analysis. Neither measures what happens to engine dependency when a
game is deliberately restructured.

\subsection{Research Gap}\label{sec:gap}

Table~\ref{tab:gap} summarizes what the surveyed literature covers and what
remains insufficiently addressed for the question of this paper.

The gap we claim is deliberately narrow. It is \emph{not} that architectural
simplification is unstudied: architectural technical debt research studies the
conditions under which structure becomes a liability
\cite{verdecchia2020,verdecchia2021,li2022}, and refactoring research studies
the removal of dispensable elements \cite{fowler2018,romano2020}. It is also
not that mixed metric outcomes are unknown: they are reported repeatedly
\cite{szoke2017,bavota2015,palomba2018}.

What we did not find covered is the combination of three properties in a single
study: (i) an intervention whose decision rule is the concrete responsibility of
each existing abstraction, applied to a working engine-based system, with
retention decisions recorded as results rather than as non-events; (ii) an
explicit control state that separates the metric effect of deleting unreachable
code from the metric effect of the architectural intervention; and (iii) an
analysis that separates within-component change from population composition
effects, so that aggregate movements are not read as changes in surviving code.

\begin{table*}[t]
\centering
\caption{What the surveyed literature covers, and what remains insufficiently
addressed for the question studied here.}
\label{tab:gap}
\small
\begin{tabular}{p{0.20\textwidth}p{0.38\textwidth}p{0.34\textwidth}}
\toprule
\textbf{Topic} & \textbf{Existing literature covers} & \textbf{Insufficiently addressed} \\
\midrule
Refactoring and metrics &
Catalogues of transformations and smells; large-scale evidence that effects on
internal quality are mixed and often small \cite{fowler2018,mens2004,szoke2017,bavota2015,palomba2018} &
Whether an observed aggregate movement reflects change in surviving components
or a change in the population being averaged \\
\addlinespace
Architecture and maintainability &
Composite and multi-property maintainability models, their calibration and
their limits \cite{coleman1994,heitlager2007,baggen2012,riaz2009,ardito2020} &
Behaviour of these indicators under an intervention that reduces the number of
components rather than increasing it \\
\addlinespace
Technical debt &
Grounded theory of architectural technical debt; erosion symptoms and their
detection \cite{verdecchia2020,verdecchia2021,kruchten2012,li2022} &
Applied interventions with before/after measurement and a control for
confounding deletions \\
\addlinespace
Game architecture and engine coupling &
Engine architecture recovery and comparison; subsystem coupling patterns;
Unity-specific smell catalogues and detectors \cite{politowski2025,ullmann2023,nardone2023,borrelli2020} &
Whether restructuring a game changes the set of components that depend on the
engine, as opposed to the proportion \\
\addlinespace
Simplification and removal of abstractions &
Dispensable smells; empirical study of dead code and unused code
\cite{fowler2018,romano2020,eder2012} &
Removal decided per abstraction by its concrete responsibility, with retention
and rejected hypotheses reported as results \\
\bottomrule
\end{tabular}
\end{table*}


%% =========================================================================
\section{Case Study Design}\label{sec:design}
%% =========================================================================

The study follows the guidelines of Runeson and H\"ost \cite{runeson2009} for
case study research in software engineering. It is a single-case, embedded,
\emph{improving} case study: the object is one system, the units of analysis
are its components, and the researchers intervened in the object and measured
the result.

\subsection{Molz and Study Context}

Molz is a mobile serious game for organic chemistry, implemented in Unity with
C\#. It represents the solubility classification procedure used to identify
unknown organic compounds: the player performs virtual tests with solvents and
reagents, observes simulated outcomes, and narrows the set of candidate
compounds by elimination. The chemical responses are driven by a structured
dataset rather than by a physical simulation, and the game offers three modes
with decreasing levels of support.

The system was developed in an interdisciplinary academic context at the
Universidade Federal do Agreste de Pernambuco, with contributions from
computing, chemistry, education and design, and is publicly available as an
open research artefact.\footnote{\url{https://github.com/bcccoworkingufape/game-quimica-2D}}
Development followed Scrum with sprints tracked on GitHub Projects, and the
architecture was documented as a layered Model-View-Presenter organization with
Data, Domain, Core and Presentation layers. A user evaluation with 103
participants was conducted and reported separately; it is context for this
study, not evidence within it.

An earlier 3D version of the game exists in the project's history. It is
\emph{not} a state of this experiment and is not used as a baseline or as a
point of comparison. The entire experiment reported here takes place inside the
2D system.

Molz is a suitable object for this question for three reasons. It is small
enough that every abstraction can be individually audited and every decision
recorded, at 91 source files and roughly 5\,000 lines of code in the baseline.
It is real: it is a released product with users, not a constructed example. And
it is engine-based, so that engine dependency is a genuine design constraint
rather than a stylistic preference.

\subsection{Experimental States C, D0 and D}

Three snapshots of \texttt{Assets/Scripts} were defined, each pinned to an
immutable commit and recovered for measurement in an isolated worktree
(Table~\ref{tab:states}).

\begin{table}[htb]
\centering
\caption{Experimental states. Each is pinned to a commit and a branch; C and D
also carry annotated tags.}
\label{tab:states}
\small
\begin{tabular}{llp{0.36\columnwidth}}
\toprule
\textbf{State} & \textbf{Commit} & \textbf{Role} \\
\midrule
C  & \texttt{9b9816d} & Pre-intervention baseline \\
D0 & \texttt{a788ea8} & Dead-code removal control \\
D  & \texttt{934eea6} & Simplified state \\
\bottomrule
\end{tabular}
\end{table}

The three states support three comparisons:

\begin{itemize}\itemsep2pt
\item $\text{C} \rightarrow \text{D0}$ --- the effect of removing dead code (control);
\item $\text{D0} \rightarrow \text{D}$ --- the effect of the architectural intervention (principal comparison);
\item $\text{C} \rightarrow \text{D}$ --- the total observed transformation.
\end{itemize}

D0 is not a third competing architecture. It exists because an audit of the
baseline found that a substantial part of the Domain layer had no consumers.
Deleting unreachable code changes every aggregate mean, because the deleted
population is not a random sample of the system. Without D0, that change would
have been indistinguishable from the effect of the architectural work.

\subsection{Architectural Simplification Protocol}

The intervention was not the application of a target architectural style. An
initial hypothesis of migrating the system to Clean Architecture / Ports and
Adapters was formulated, audited against the baseline, and \emph{rejected before
any code was written}: state C already had a Unity-free domain layer, separated
Data, Domain and Presentation folders, and abstractions carrying concrete
responsibilities. The decision rule adopted instead was:

\begin{quote}
\emph{For each existing abstraction, identify the responsibility it performs in
the running system. Retain the abstraction if and only if that responsibility is
concrete and is not already discharged elsewhere.}
\end{quote}

Four outcomes were possible for each abstraction: REMOVE, MERGE, STRENGTHEN
(the abstraction is justified but bypassed in practice) and KEEP. KEEP decisions
were recorded with their evidence, on the same footing as changes, so that a
retention is reportable as a result rather than as the absence of one.

The intervention proceeded in eight stages (Table~\ref{tab:stages}). Each stage
was implemented in isolation, compiled, verified against the test suite, and
approved by a human reviewer before the next stage began. Stage A alone
produced D0; stages B through H produced D.

\begin{table}[htb]
\centering
\caption{Intervention stages. Stage A produced the control state D0; stages B--H
produced the simplified state D.}
\label{tab:stages}
\small
\begin{tabular}{cp{0.68\columnwidth}}
\toprule
\textbf{Stage} & \textbf{Intervention} \\
\midrule
A & Dead-code removal: 11 files, 20 types \\
B & Data layer: one DTO and four mappers removed; transformation localized in the repositories \\
C & Audio boundary adopted: 56 concrete accesses to engine singletons reduced to 0 \\
D & Partial explicit composition: three dead registrations removed; one presenter composed by constructor; service registry deliberately retained \\
E & Score multiplier consolidated into the domain scoring service \\
F & Indicator-reagent semantics given a single owner in the domain entity \\
G & View\,$\leftrightarrow$\,orchestrator cycles broken through two role contracts \\
H & Panel-reset duplication replaced by the existing collaborator \\
\bottomrule
\end{tabular}
\end{table}

Two constraints were imposed on the intervention and observed throughout. No
\texttt{MonoBehaviour} was renamed, moved between namespaces, or relocated in a
way that would invalidate an asset GUID, because that would silently break
scene serialization. And no global access mechanism was to be replaced by
another global access mechanism: a change that substitutes one ambient
singleton for another was not counted as simplification.

\subsection{Functional-Parity Protocol}

The criterion for the experiment is $\mathrm{Behaviour}(\text{C}) =
\mathrm{Behaviour}(\text{D})$. Because the object is an engine-based game whose
wiring lives partly in serialized scene data, no single verification technique
is sufficient. Five were combined.

\emph{Characterization and invariant tests.} A suite of 118 tests was written
against state C before any modification and kept unchanged thereafter. It
covers data loading and deserialization, solubility rules, gameplay rules,
session state, entity invariants, repository materialization, scoring, indicator
semantics, and two architecture-boundary invariants enforced by source scanning.
The suite compiles the engine-free production files outside Unity; it is not
part of the population measured by the instrument.

The suite deliberately preserves several behaviours that would ordinarily be
treated as defects, including a difficulty asset with a zero score multiplier,
a divergence in the default difficulty, an inert branch in a UI controller, and
the specific parsing policy applied to malformed dataset entries. Correcting any
of them would have changed behaviour and broken the parity criterion.

\emph{Compilation.} Both through an isolated project definition and through the
engine's own assembly, in each state.

\emph{Engine import.} The Unity project was imported in state D with zero
missing-script errors and zero serialization errors, confirming that no GUID or
serialized binding was invalidated.

\emph{Manual full-flow verification.} A scripted walkthrough of the complete
game flow in all three modes.

\emph{Smoke verification} after the final stage, covering phase restart, panel
reset and return to menu.

\subsection{Measurement Instrument}

All three states were measured with a single instrument, frozen at one commit
and never read from inside the measured worktrees. It consists of a Roslyn-based
metrics extractor built on \texttt{Microsoft.CodeAnalysis.CSharp} 4.11.0, plus
SonarAnalyzer.CSharp 10.31.0.145097 and the engine's own analyzers, executed on
.NET SDK 9.0.306 under Unity 6000.4.1f1. Each state is recovered with a detached
worktree at its pinned commit, so production content and measurement tooling
remain separate; no snapshot is modified to accommodate the instrument.

The instrument reports, per type, lines of code, efferent coupling, and whether
the type references \texttt{UnityEngine}; and, per member, cyclomatic complexity
and maintainability index. Aggregations are computed from the raw per-type and
per-member records rather than from summary output.

Two properties of the instrument were verified explicitly. \emph{Determinism}:
re-executing the pipeline on state D produced byte-identical artefacts,
including the analyzer SARIF. \emph{Version equivalence}: an error in the
project-generation step was found and corrected during the study; after the
correction, state C was re-measured through the real entry point and the
resulting files were byte-identical to those produced before, so all three
states are measured by the same instrument version.

Analyzer violation counts were also collected (C 197, D0 191, D 143). They are
auxiliary pipeline output and are not used to answer any research question here;
interpreting individual findings belongs to a separate study.

\subsection{Component Classification}

To answer RQ3, every source file was assigned to exactly one category. The
assignment was derived mechanically from the commit history --- which files each
stage touched --- rather than from post-hoc judgement.

\begin{description}\itemsep2pt
\item[DEAD CODE REMOVED] Present in C, removed exclusively in $\text{C}\rightarrow\text{D0}$,
after triple evidence of unreachability: no C\# references, no GUID references in
scenes, prefabs or assets, and no occurrence of the type name in serialized YAML.
\item[DIRECTLY REFACTORED] Component whose code or responsibility was a deliberate
target of stages B--H.
\item[INDIRECTLY AFFECTED] Component changed only to accommodate a change whose
primary target was another component.
\item[PRESERVED] Existing component deliberately left unmodified by the relevant
architectural interventions.
\item[NEW INTERVENTION ARTIFACT] Element absent from the baseline, introduced as a
necessary part of the intervention.
\end{description}

A component that does not exist in a given state is recorded as \emph{absent},
never as zero. A NEW INTERVENTION ARTIFACT therefore has no row for C or D0, and
DEAD CODE REMOVED has no row for D0 or D. Substituting zero would inject
artificial values into every mean.

\subsection{Analysis Procedure}

Three analyses were performed on the frozen measurements.

\emph{Aggregate comparison.} All summary values were recomputed from the raw
per-type and per-member records in each state and cross-checked against the
instrument's own summary; all 42 values matched.

\emph{Matched-population analysis.} Aggregate means confound two phenomena:
change inside surviving components, and change in the population being averaged.
Components were therefore matched across states by a stable identity key ---
\texttt{file|type} for types and \texttt{file|type|member} for members --- and
the mean was recomputed over the intersection only. Keys that were not unique
within a state, which occurred for two method overloads in each state, were
\emph{excluded} from matching rather than resolved by assumption.

\emph{Category analysis.} Metrics were aggregated within each RQ3 category, with
units kept explicit: lines of code, efferent coupling and type counts are
per-type quantities; cyclomatic complexity and maintainability index are
per-member quantities, and group-level values for them are means over the
members belonging to the group, never aggregated per type.

No statistical significance testing was performed. The study measures one
system at three points in time; there is no sample over which a test would be
defined, and no test is reported.


%% =========================================================================
\section{Results}\label{sec:results}
%% =========================================================================

\subsection{RQ1 --- Architectural Reorganization}\label{sec:rq1}

Table~\ref{tab:rq1} records the outcome of the responsibility assessment for
each abstraction examined. Thirteen concerns were audited: three led to removal or
merging, two to strengthening, three to a local simplification or a change
of ownership, and five to a documented retention.

\begin{table*}[t]
\centering
\caption{RQ1 --- responsibility assessment of each abstraction. Retention
decisions are reported with the same status as changes.}
\label{tab:rq1}
\small
\begin{tabular}{p{0.16\textwidth}p{0.13\textwidth}p{0.32\textwidth}p{0.13\textwidth}p{0.16\textwidth}}
\toprule
\textbf{Concern} & \textbf{Initial hypothesis} & \textbf{Evidence found} & \textbf{Decision} & \textbf{Consequence} \\
\midrule
Question mapper & Redundant & Purely positional pass-through; invariants live in the entity & REMOVE & Binding now rests on constructor parameter names \\
\addlinespace
Question DTO & Redundant & 1:1 copy, no type conversion & REMOVE & Verified by a nine-criterion spike before removal \\
\addlinespace
Three further mappers & Duplicated & One caller each; real conversion present & MERGE & Transformation became a private repository method \\
\addlinespace
Remaining DTOs & Useful & Case conversion, string-to-enum, one rename & KEEP & Representation boundary preserved \\
\addlinespace
Repository surface & Cleanup candidate & Four methods without consumers, three of which lost them \emph{in Stage A} & KEEP & Removing them would credit the intervention with a control effect \\
\addlinespace
Audio port & Necessary but bypassed & 56 concrete accesses against 2 port consumers & STRENGTHEN & 56 $\rightarrow$ 0; port gained two availability members \\
\addlinespace
Service locator & \textbf{Redundant} & Full removal would require scene scanning, layer inversion or a lifetime change & \textbf{HYPOTHESIS REJECTED} $\rightarrow$ KEEP + SIMPLIFY & Three dead registrations removed; registry retained for cross-scene lifetime \\
\addlinespace
Game manager responsibilities & Broad decomposition & Only scoring was split across layers & Selective SIMPLIFY & One move, seven documented retentions \\
\addlinespace
Indicator identity & Duplicated in Presentation & Confirmed at two sites & MOVE & Ownership moved to the domain entity \\
\addlinespace
Solvent group buttons & Duplicated & The consumer sustaining the duplication was removed \emph{in Stage A} & \textbf{HYPOTHESIS REJECTED} $\rightarrow$ KEEP & Content is editorial and carries markup \\
\addlinespace
View\,$\leftrightarrow$\,orchestrator cycles & Breakable & Two concrete reverse dependencies, two engine-type leaks & STRENGTHEN & Two role contracts introduced; the existing view contract untouched \\
\addlinespace
God-object candidates & Decomposition & Only one duplication had a better existing owner & Local SIMPLIFY & One consolidation, nine retentions \\
\addlinespace
Pre-existing dead members & Dead & Already dead in C & PRESERVE FOR PARITY & Retained so control and intervention stay separable \\
\bottomrule
\end{tabular}
\end{table*}

Structurally, the intervention produced a system with \emph{fewer classes and
more interfaces}: from D0 to D the class count fell from 65 to 60, the interface
count rose from 12 to 14, and the member count rose from 451 to 457. Five files
were removed and two were added. This is the opposite of the shape usually
associated with simplification, and it follows directly from the decision rule:
abstractions were removed where they carried no responsibility, and introduced
where a responsibility existed that no abstraction expressed --- specifically,
two role contracts that broke concrete cycles between views and orchestrators.

\medskip\noindent\textbf{Four mechanisms of abstraction utility.} The
usefulness of an abstraction turned out to be context-dependent rather than
determined by the pattern it resembled. Four distinct mechanisms appeared in the
audit, and they are epistemologically different:

\begin{enumerate}\itemsep2pt
\item \emph{Abstraction without a current responsibility} --- removed. The question
mapper performed a positional copy that the entity constructor already performed.
\item \emph{Abstraction with a concrete responsibility but bypassed in practice} ---
strengthened. The audio port was well-formed, but 56 call sites reached the engine
singletons directly instead of going through it.
\item \emph{Abstraction whose responsibility only the removal attempt revealed} ---
retained. The service locator appeared redundant under static inspection; the
attempt to remove it showed that it carried a cross-scene lifetime responsibility
that no other mechanism in the system discharged. Here static inspection was
\emph{incomplete}.
\item \emph{Abstraction whose justification an earlier transformation had already
removed} --- retained. The duplication in the solvent group buttons had a consumer
when it was first identified, and that consumer was deleted in Stage A. Here
static inspection was \emph{correct at the time} and became obsolete.
\end{enumerate}

Cases 3 and 4 both produce a retained abstraction, but for opposite reasons: in
one the analysis was wrong, in the other the system moved underneath a correct
analysis. Reporting them as a single category of ``false positives'' would lose
the distinction.

\subsection{RQ2 --- Internal Quality Indicators}\label{sec:rq2}

\subsubsection{Aggregate measurements}

Table~\ref{tab:aggregate} reports the full aggregate profile of the three
states.

\begin{table}[htb]
\centering
\caption{Aggregate measurements across the three states. CC and MI are
per-member quantities; LOC, coupling and type counts are per-type.}
\label{tab:aggregate}
\small
\begin{tabular}{lrrr}
\toprule
\textbf{Indicator} & \textbf{C} & \textbf{D0} & \textbf{D} \\
\midrule
Source files            & 91    & 80    & 77 \\
LOC (non-blank)         & 5\,093 & 4\,738 & 4\,729 \\
Comment lines           & 695   & 649   & 825 \\
Classes                 & 79    & 65    & 60 \\
Interfaces              & 16    & 12    & 14 \\
Structs / records       & 7     & 5     & 5 \\
Enums                   & 9     & 9     & 9 \\
Members                 & 496   & 451   & 457 \\
\addlinespace
CC mean                 & 2.33  & 2.39  & 2.38 \\
CC median               & 2.0   & 2.0   & 2.0 \\
CC max                  & 17    & 17    & 17 \\
CC total                & 1\,158 & 1\,078 & 1\,089 \\
\addlinespace
MI mean                 & 72.5  & 71.7  & 71.8 \\
MI median               & 71.2  & 70.2  & 70.9 \\
Members with MI $\geq$ 65 & 323 & 287   & 293 \\
\addlinespace
Efferent coupling mean  & 11.4  & 12.8  & 13.4 \\
Engine-coupled classes  & 24    & 23    & 23 \\
\quad as \% of classes  & 30.4  & 35.4  & 38.3 \\
\bottomrule
\end{tabular}
\end{table}

Between D0 and D the system lost 3 files, 9 lines of code and 5 classes, and
gained 2 interfaces and 6 members. The three principal indicators moved very
little: MI mean rose by 0.1 on a 0--100 scale, CC mean fell by 0.01 on a scale
whose observed mean is about 2.4, and mean efferent coupling rose by 0.6.

Between C and D0 --- the control --- MI mean fell by 0.8 and mean efferent
coupling rose by 1.4, even though the only action taken was the deletion of
unreachable code.

\subsubsection{Matched-population analysis}

The aggregate figures above mix two phenomena: change inside components that
exist in both states, and change in the population over which the mean is taken.
Table~\ref{tab:matched} separates them.

\begin{table}[htb]
\centering
\caption{Matched-population analysis, D0 $\rightarrow$ D. The sign of the
system-wide MI and CC movement inverts once the population is held fixed.}
\label{tab:matched}
\small
\begin{tabular}{lrrr}
\toprule
\textbf{Basis} & \textbf{D0} & \textbf{D} & \textbf{$\Delta$} \\
\midrule
\multicolumn{4}{l}{\emph{Members, system-wide (n = 442 matched)}} \\
MI, aggregate mean      & 71.7  & 71.8  & $+0.1$ \\
MI, matched members     & 71.87 & 71.80 & $-0.06$ \\
CC, aggregate mean      & 2.39  & 2.38  & $-0.01$ \\
CC, matched members     & 2.41  & 2.42  & $+0.01$ \\
\addlinespace
\multicolumn{4}{l}{\emph{Directly refactored types}} \\
Coupling, matched (n = 19) & 26.84 & 28.37 & $+1.53$ \\
Coupling, raw group mean   & 22.29 & 28.37 & $+6.08$ \\
Coupling of removed types  & 5.00  & ---   & --- \\
\bottomrule
\end{tabular}
\end{table}

Three observations follow.

\emph{First, the signs invert.} The aggregate MI mean rises while matched-member
MI falls; the aggregate CC mean falls while matched-member CC rises. Both
aggregate movements are composition effects, driven by the eleven new members ---
largely simple properties with high MI --- and the five removed ones. Reporting
the aggregate delta alone would have supported the statement that maintainability
improved slightly; the matched analysis shows that within surviving code it did
not.

\emph{Second, the magnitudes are negligible either way.} Only 65 of 442 matched
members (14.7\%) changed MI at all, and only 12 (2.7\%) changed CC. On scales of
0--100 and roughly 2.4, movements of 0.06 and 0.01 carry no practical weight in
either direction.

\emph{Third, the coupling increase is real and partly genuine.} Of the $+6.08$
raw increase in the directly refactored group, $+1.53$ is within-type change
among the 19 surviving types, and the remainder follows from removing five types
whose mean coupling (5.00) was far below the group mean. Fourteen of the 19
survivors changed, twelve upwards and two downwards; the largest increases were
in the repositories that absorbed the mapper logic --- one rose from 12 to 19,
another from 14 to 19, another from 12 to 15. Architectural responsibility
redistribution and measured coupling increase are related here, but they are not
the same phenomenon, and the increase is not merely a redistribution artefact.

\subsubsection{Engine dependency}

The set of types referencing \texttt{UnityEngine} was compared by identity, not
by count:

\begin{center}\small
\begin{tabular}{lr}
Engine-coupled types in D0 & 23 \\
Engine-coupled types in D  & 23 \\
Intersection               & 23 \\
Present only in D0         & --- (empty) \\
Present only in D          & --- (empty) \\
\end{tabular}
\end{center}

The set is identical. The percentage rose from 35.4\% to 38.3\% because the
class population fell from 65 to 60 while the absolute count stayed constant.
The movement of that percentage is therefore a denominator effect and carries no
information about engine dependency.

This is the clearest single result of the study, and it contradicted an
expectation that had been central to the earlier framing of the work. The
components that depend on the engine are the components the engine instantiates
and wires --- the \texttt{MonoBehaviour} classes that scenes reference by GUID.
Restructuring the layers above them did not change which components those are.

\subsection{RQ3 --- Distribution Across Component Categories}\label{sec:rq3}

Table~\ref{tab:rq3dist} gives the distribution of files across the five
categories in each state, and Table~\ref{tab:rq3metrics} the group metrics.

\begin{table}[htb]
\centering
\caption{RQ3 --- distribution of source files by category. A dash denotes
absence, not zero.}
\label{tab:rq3dist}
\small
\begin{tabular}{lrrr}
\toprule
\textbf{Category} & \textbf{C} & \textbf{D0} & \textbf{D} \\
\midrule
Directly refactored        & 24 & 24 & 19 \\
Indirectly affected        & 1  & 1  & 1 \\
Preserved                  & 55 & 55 & 55 \\
Dead code removed          & 11 & ---& --- \\
New intervention artifact  & ---& ---& 2 \\
\midrule
\textbf{Total}             & \textbf{91} & \textbf{80} & \textbf{77} \\
\bottomrule
\end{tabular}
\end{table}

\begin{table*}[t]
\centering
\caption{RQ3 --- group metrics. The upper block is the intervention
(D0 $\rightarrow$ D); the lower block is the control (C $\rightarrow$ D0).}
\label{tab:rq3metrics}
\small
\begin{tabular}{lrrrrrrr}
\toprule
\textbf{Group} & \textbf{Types} & \textbf{Members} & \textbf{LOC} & \textbf{CC mean} & \textbf{MI mean} & \textbf{Coupling mean} & \textbf{Engine} \\
\midrule
Directly refactored, D0   & 24 & 199 & 2\,357 & 2.44 & 71.2 & 22.29 & 8 \\
Directly refactored, D    & 19 & 204 & 2\,505 & 2.43 & 71.4 & 28.37 & 8 \\
\addlinespace
Indirectly affected, D0   & 3  & 11  & 58    & 1.36 & 83.5 & 7.00  & 0 \\
Indirectly affected, D    & 3  & 12  & 67    & 1.33 & 84.0 & 7.67  & 0 \\
\addlinespace
\textbf{Preserved, D0}    & \textbf{64} & \textbf{241} & \textbf{3\,286} & \textbf{2.40} & \textbf{71.5} & \textbf{9.56} & \textbf{15} \\
\textbf{Preserved, D}     & \textbf{64} & \textbf{241} & \textbf{3\,286} & \textbf{2.40} & \textbf{71.5} & \textbf{9.56} & \textbf{15} \\
\addlinespace
New intervention artifact, D & 2 & 0 & 30 & n/a & n/a & 0.50 & 0 \\
\midrule
\multicolumn{8}{l}{\emph{Control, C $\rightarrow$ D0}} \\
Dead code removed (C only) & 20 & 45  & 356 & \textbf{1.78} & \textbf{80.3} & \textbf{4.50} & --- \\
Whole project in C, for contrast & 111 & 496 & --- & 2.33 & 72.5 & 11.4 & --- \\
\bottomrule
\end{tabular}
\end{table*}

The new intervention artifacts have zero members because both are pure
interfaces; the instrument assigns no cyclomatic complexity or maintainability
index to declarations without bodies. Their coupling of 0.50 is reported as
measured.

\medskip\noindent\textbf{The preserved group is metric-identical.} Across
D0~$\rightarrow$~D, the 64 preserved types show no change in any reported
indicator: same type count, same member count, same lines of code, same
complexity, same maintainability, same coupling, same engine-coupled count. In
the remaining groups the control transition C~$\rightarrow$~D0 is likewise near
zero: $-7$ lines and $-0.12$ coupling in the directly refactored group, $-0.02$
coupling in the preserved group, and no change at all in the indirectly affected
group.

\medskip\noindent\textbf{The removed dead code was not representative.} The 20
types deleted in the control had mean MI 80.3 against 72.5 for the project, mean
cyclomatic complexity 1.78 against 2.33, and mean efferent coupling 4.50 against
11.4. They were simpler, better-scoring and less coupled than the system
average --- which is unsurprising, since code that nothing calls tends to be
small and to depend on little. Removing them therefore moved every normalized
mean unfavourably. This is the confound the control state was introduced to
isolate, and it is quantitatively large enough to have swamped the effect of the
architectural intervention had the two been measured together.

\medskip\noindent\textbf{Interpretation and its limit.} The measured structural
changes were confined to components classified as directly or indirectly
affected, while the preserved group remained metric-identical. This is evidence
that the changes were \emph{contained}. It is not an independent validation of
the classification and not a causal attribution: the preserved group is
\emph{defined} by the absence of relevant source changes, and the metrics are
derived from static code, so its metric identity partly follows from its
definition. Classification and result are not independent, and we do not claim
they are.

\subsection{Summary of Unexpected Findings}

Six results ran against the expectations held at the start of the study.

\begin{enumerate}\itemsep2pt
\item Removing dead code moved every normalized mean unfavourably, because the
removed population scored better than the system average.
\item Engine dependency did not fall: the same 23 types in D0 and D, with the
percentage rising through denominator change alone.
\item Simplification produced fewer classes but \emph{more} interfaces and more
members.
\item Two a-priori simplification hypotheses were rejected by implementation
evidence, for two different reasons.
\item The MI and CC movements invert sign under matched-population analysis.
\item Coupling rose in the repositories that absorbed the mapper logic, both by
redistribution and as a genuine within-type increase.
\end{enumerate}


%% =========================================================================
\section{Discussion}\label{sec:discussion}
%% =========================================================================

\subsection{Aggregate Indicators and Population Change}

The clearest methodological lesson of this study is that an aggregate mean over
a population that is itself being modified is a poor instrument for evaluating
an intervention that modifies the population.

Two of our three principal indicators change sign between the aggregate and the
matched view (Table~\ref{tab:matched}). Had we reported only the aggregate, the
natural narrative would have been that simplification produced a small
improvement in maintainability and a small reduction in complexity. The matched
analysis shows that within the 442 members present in both states, maintainability
moved down by 0.06 and complexity up by 0.01 --- and that the aggregate movement
came from eleven newly introduced members, mostly simple properties that score
well on MI.

Both readings are arithmetically correct. They answer different questions. The
aggregate answers ``what is the profile of the system that now exists''; the
matched analysis answers ``what happened to the code that survived''. For
evaluating an intervention, the second is the relevant one, and it is the one
the literature reports less often.

The same point applies with more force to the control transition. Removing
unreachable code is uncontroversially good practice \cite{romano2020,eder2012},
and it made every normalized indicator in this system worse. Nothing about the
surviving code changed; the composition of the population did. A study that
combined dead-code removal with an architectural intervention and reported the
aggregate delta would have attributed this movement to the architecture.

This is consistent with, and gives a concrete mechanism for, the repeatedly
reported finding that refactoring produces mixed rather than uniform metric
outcomes \cite{szoke2017,bavota2015,palomba2018}. Part of the mixture may be
population effects rather than disagreement about the value of the
transformations.

\subsection{Engine Dependency Is Structural, Not Architectural}

The engine-coupled set was identical in D0 and D: the same 23 types, with empty
differences in both directions. This is the result that most directly
contradicted the study's initial framing.

The explanation is specific to engine-based systems and consistent with what the
engine architecture literature reports about coupling to engine cores
\cite{politowski2025,ullmann2023}. In Unity, a component that a scene
instantiates must be a \texttt{MonoBehaviour}; a field that the editor
serializes must be visible to the engine's serializer; a lifetime that spans
scene loads must be managed by an engine mechanism \cite{unitydocs}. These are
not layering choices that an architect can revise. They are the terms on which the engine
composes the application.

The architectural layers above that boundary can be reorganized freely, and were.
But reorganizing them does not migrate a type across the boundary, because the
boundary is defined by what the engine must instantiate, not by what the design
would prefer. Interpreting the movement from 35.4\% to 38.3\% as increased
engine dependency would be a denominator artefact; interpreting a hypothetical
movement downwards as reduced dependency would have been equally unfounded.

The practical implication for engine-based systems is that engine dependency
should be reported as a set, and its change as a set difference. A percentage
whose denominator is the class population will move whenever the population
moves, for reasons that have nothing to do with the engine.

\subsection{Responsibility as a Decision Rule}

The intervention did not apply a target architectural style, and the a-priori
hypothesis of migrating to Clean Architecture was rejected during the audit,
before any code was written, because the baseline already had a Unity-free
domain and abstractions carrying real responsibilities.

What replaced it --- assessing each abstraction against the responsibility it
concretely performs --- produced an outcome that no style-driven rule would have
produced. It removed a DTO and four mappers, but retained three other DTOs that
performed genuine representation conversion. It retained a service locator that
a style-driven rule would have eliminated, because the removal attempt revealed a
cross-scene lifetime responsibility that nothing else discharged. And it
\emph{introduced} two interfaces, because two concrete cycles existed for which
no abstraction expressed the responsibility.

The four mechanisms identified in Section~\ref{sec:rq1} suggest a practical
ordering. Mechanisms 1 and 2 --- no responsibility, and bypassed responsibility
--- are decidable by static inspection. Mechanisms 3 and 4 are not. Mechanism 3
(the responsibility that only the removal attempt reveals) means that static
inspection can be \emph{incomplete}: some responsibilities are visible only in
the runtime composition model, which for an engine-based system lives partly in
serialized data. Mechanism 4 (the justification that an earlier transformation
already removed) means that a correct analysis can become obsolete during a
staged intervention, which is an argument for re-deriving candidates after each
stage rather than fixing the plan at the outset.

This connects to the architectural technical debt literature
\cite{verdecchia2020,verdecchia2021}, which characterizes such debt as decisions
that were appropriate when taken and later became obstacles. Mechanism 4 is the
inverse-and-benign case: an item of apparent debt whose justification
disappeared, making the item harmless rather than harmful, and detectable only
by re-examination.

\subsection{What Simplification Cost and What It Bought}

The intervention did not buy a better metric profile. Within surviving code the
indicators were flat to marginally worse, and mean efferent coupling rose both
in aggregate ($+0.6$) and among matched surviving types in the directly
refactored group ($+1.53$).

The coupling increase has a concrete cause. Removing four mappers moved their
transformation logic into the repositories that had called them, so those
repositories now depend on more types than before. Whether this is an
improvement is a design judgement, not a metric one: the system has fewer
indirection layers and fewer components to keep consistent, at the cost of
repositories that know more. The metric records the cost and is silent about the
benefit --- which is precisely the limitation that the multi-property
maintainability models were designed to address \cite{heitlager2007,baggen2012},
and a reason not to treat a coupling mean as a verdict.

What the intervention did produce is structural: 56 direct accesses to engine
audio singletons reduced to zero behind an existing port; two concrete cycles
between views and orchestrators replaced by role contracts; one duplicated
domain rule given a single owner; three dead service registrations removed. These
are changes in the dependency structure that the aggregate indicators do not
express, and reporting them requires the architectural narrative of
Section~\ref{sec:rq1} alongside the numbers.

We deliberately do not claim that state D is better than state D0. The study
measured what changed; whether the trade is worthwhile depends on which future
changes the system will actually face, and that is not observable from a static
snapshot.

\subsection{Implications}

For \emph{researchers}: interventions that change component population should
report matched-population results alongside aggregate ones, and should separate
deletion effects from restructuring effects. In this study the deletion effect
was larger than the restructuring effect on every aggregate indicator.

For \emph{practitioners} working on engine-based systems: engine dependency is
best tracked as a set of types rather than as a percentage, and abstractions
should be assessed against the responsibility they discharge in the composed,
running system --- not against the catalogued pattern they resemble in source
\cite{gamma1995}.

For both: a retention decision, recorded with its evidence, is a result. Five of the thirteen audited
concerns in this study ended in a retention decision, and two of those
overturned an a-priori hypothesis. A protocol that only records changes would have reported
none of it.


%% =========================================================================
\section{Threats to Validity}\label{sec:threats}
%% =========================================================================

We follow the categories of Runeson and H\"ost \cite{runeson2009}.

\medskip\noindent\textbf{Construct validity.} MI, cyclomatic complexity and
efferent coupling are structural proxies, not measurements of the maintainability
that developers experience. The critiques are well established: cyclomatic
complexity is substantially a proxy for size \cite{shepperd1988}, and MI is a
composite calibrated on a corpus unlike ours that incorporates both
\cite{coleman1994,ardito2020}. We mitigate this by reporting the indicators
separately as well as together, by reporting medians alongside means, and by
never treating a movement in one of them as a statement about maintainability.
For the same reason we use the term \emph{internal software quality indicators}
throughout. \emph{Engine dependency} is defined operationally as a count of types
referencing \texttt{UnityEngine} as determined by the instrument; its percentage
uses classes as denominator and is therefore sensitive to population change, which
is why we report the set identity rather than the percentage as the result.

\medskip\noindent\textbf{Internal validity.} The eight stages were applied
sequentially and may interact. The D0 checkpoint isolates dead-code removal from
the architectural intervention, but it does not isolate stages B through H from
one another. Two concrete interactions were in fact observed and are reported in
Table~\ref{tab:rq1}: in both, Stage A removed the justification for a change
that later stages had been expected to make. We did not measure after every
stage, so finer attribution within B--H is not available.

\medskip\noindent\textbf{External validity.} This is a single case: one
Unity serious game of 77 files in its final state, developed by a small academic
team, with three sequentially loaded scenes. The conclusions about cross-scene
lifetime depend on that scene topology. The conclusions about engine dependency
should generalize to engine-based systems with a comparable composition model,
but that is an expectation, not a demonstration. Nothing here supports
generalization to large systems, to other engines, or to teams with different
practices.

\medskip\noindent\textbf{Reliability.} The instrument is versioned and frozen at
a single commit; each state is pinned by commit, tag and branch and recovered in
an isolated worktree; determinism was confirmed by a byte-identical re-run of the
whole pipeline including the analyzer output. An instrument defect found during
the study was corrected, after which state C was re-measured and produced
byte-identical results, establishing that all three states were measured by the
same version. The measured system, the instrument, the raw per-type and
per-member records and the analysis scripts are in the public repository.

\medskip\noindent\textbf{Researcher bias.} The authors performed both the
intervention and its analysis, which is the standard threat in improving case
studies. Partial mitigations: the classification criteria for RQ3 were fixed
before measurement; the RQ3 assignment is derived mechanically from the commit
history rather than from judgement after seeing the numbers; each stage passed a
human approval gate before the next began; and a list of prohibited claims ---
including several that the data superficially appeared to support --- was fixed
before the manuscript was written. We did not seek to demonstrate that D is
better than D0, and we report the results that contradicted our initial framing.

\medskip\noindent\textbf{Functional equivalence.} The parity evidence is strong
but it is evidence, not proof. The 118 characterization and invariant tests, the
clean engine import, the manual full-flow walkthrough and the post-stage smoke
verification cover the behaviour they cover. We claim behavioural preservation
for that scope and no more; we do not claim formal functional equivalence.
Structural differences on paths that no verification reaches were recorded rather
than hidden.

\medskip\noindent\textbf{LLM-assisted implementation.} The implementation was
carried out with the assistance of an LLM coding agent, under a protocol with
per-stage human approval and explicit prohibitions. This may have influenced the
shape of the code and the set of options considered at each stage. The study
neither measures nor claims anything about AI-assisted productivity or quality;
no aspect of the design would support such a claim.


%% =========================================================================
\section{Conclusion}\label{sec:conclusion}
%% =========================================================================

This paper reported a comparative case study of responsibility-driven
architectural simplification in Molz, a Unity-based serious game, measured across
three pinned states of the same system with a single frozen instrument: the
baseline C, a dead-code-removal control D0, and the simplified state D.

Answering RQ1, the reassessment of thirteen architectural concerns produced five
removals or merges, three strengthenings and five documented retentions, two of
which overturned an a-priori hypothesis. The resulting system has fewer classes
(65 to 60) but more interfaces (12 to 14) and more members (451 to 457):
simplification here meant removing abstractions that carried no responsibility
while introducing two that expressed one.

Answering RQ2, the aggregate indicators moved very little between D0 and D --- MI
71.7 to 71.8, CC 2.39 to 2.38, mean efferent coupling 12.8 to 13.4 --- and a
matched-population analysis over the 442 members present in both states inverted
the sign of the first two, showing the aggregate movements to be composition
effects rather than changes in surviving code. The set of engine-coupled types
was identical in both states; only its percentage moved, through denominator
change.

Answering RQ3, the measured changes were contained within the directly and
indirectly affected groups, while the 64 preserved types remained
metric-identical. The dead code removed in the control had a substantially better
metric profile than the system average (MI 80.3 against 72.5; coupling 4.50
against 11.4), so its removal moved every normalized mean unfavourably --- an
effect that, without the control state, would have been credited to the
architectural intervention.

The study's central methodological claim is that interventions which change the
population of components cannot be evaluated by aggregate means alone. Its
central empirical observation is that in an engine-based system, the set of
components depending on the engine is fixed by the engine's composition model
and is not a variable that layer reorganization can move.

Several directions remain open. Measuring after each stage rather than only at
the endpoints would permit attribution within the intervention. Repeating the
design on systems built with other engines, and on systems large enough for the
population effects to be studied statistically, would test how far the
observations extend. A companion study of the analyzer findings collected here,
which this paper deliberately leaves uninterpreted, is a separate line of work.
Finally, the responsibility-driven decision rule was applied by the authors; how
consistently different engineers would apply it to the same system is an
empirical question we did not address.

\begin{acknowledgements}
%% NEEDS SOURCE: funding programme identifiers and the definitive list of
%% acknowledged contributors are to be confirmed before submission.
This work was carried out within the PIBIC 2025--2026 undergraduate research
programme at the Universidade Federal do Agreste de Pernambuco. The authors
thank the members of the Molz project team for their contributions to the
system studied here.
\end{acknowledgements}

\bibliographystyle{apalike}
\bibliography{paper1-references}

\end{document}
